% =============================================================================
% FICHAMENTO DE LEITURA — MODELO ESTRUTURADO
% =============================================================================
% Características do fichamento:
%   - Registro sistemático de leitura
%   - Organizado por autor/obra
%   - Combina citações diretas, paráfrases e análise crítica
%   - Útil para revisão de literatura, dissertações e teses
% =============================================================================
\documentclass[
  12pt,
  a4paper,
  brazilian
]{article}

% --- Pacotes ---
\usepackage[utf8]{inputenc}
\usepackage[T1]{fontenc}
\usepackage[brazilian,portuges]{babel}
\usepackage[alf]{abntex2cite}
\usepackage{abntex2abrev}
\usepackage{setspace}
\usepackage{geometry}
\usepackage{hyperref}
\usepackage{xcolor}
\usepackage{enumitem}
\usepackage{longtable}
\usepackage{booktabs}
\usepackage{microtype}

% --- Margens ---
\geometry{
  a4paper,
  left=3cm,
  top=3cm,
  right=2cm,
  bottom=2cm
}

\onehalfspacing

\hypersetup{
  colorlinks=true,
  linkcolor=black,
  citecolor=black,
  urlcolor=black
}

\title{Fichamentos de Leitura}
\author{Nome do Pesquisador}
\date{2026}

\begin{document}

\maketitle

\section*{Instruções}

Cada fichamento deve conter:
\begin{enumerate}
  \item \textbf{Referência completa} da obra (ABNT NBR 6023:2018)
  \item \textbf{Palavras-chave} representativas do conteúdo fichado
  \item \textbf{Citações diretas} transcritas com página
  \item \textbf{Paráfrases} com indicação de página
  \item \textbf{Comentário crítico} do pesquisador
  \item \textbf{Conexão} com outros autores e com a própria pesquisa
\end{enumerate}

\newpage

% ============================================================================
% FICHAMENTO 1
% ============================================================================
\section*{Fichamento 1}
\addcontentsline{toc}{section}{Fichamento 1 — Autor, Ano}

\begin{longtable}{p{3cm} p{12cm}}
\toprule
\textbf{Obra} & SOBRENOME, Nome. \textit{Título do Livro}: subtítulo.
               Edição. Cidade: Editora, Ano. \\
\textbf{Palavras-chave} & Conceito A; Conceito B; Metodologia C \\
\midrule
\textbf{Citações diretas} &
  ``A citação direta com indicação de página (Autor, ano, p. X).'' \\
&
  ``Outra citação direta relevante (Autor, ano, p. Y).'' \\
\midrule
\textbf{Paráfrases} &
  No parágrafo que trata de [tema], o autor argumenta que... (Autor, ano, p. Z). \\
&
  O conceito de [conceito] é definido como... (Autor, ano, p. W). \\
\midrule
\textbf{Comentário crítico} &
  \textit{Análise do pesquisador:} este conceito dialoga diretamente
  com o referencial teórico adotado na dissertação, especialmente com
  a noção de [conceito relacionado] desenvolvida por [outro autor]. \\
&
  \textit{Questão em aberto:} como operacionalizar este conceito em
  uma pesquisa empírica com [características do contexto]? \\
\midrule
\textbf{Conexões} &
  \begin{itemize}[nosep]
    \item Dialogar com: \cite{autor2023}
    \item Contradizer: \cite{autor2021}
    \item Aprofundar: \cite{autor2020}
  \end{itemize} \\
\bottomrule
\end{longtable}

\newpage

% ============================================================================
% FICHAMENTO 2
% ============================================================================
\section*{Fichamento 2}
\addcontentsline{toc}{section}{Fichamento 2 — Autor, Ano}

\begin{longtable}{p{3cm} p{12cm}}
\toprule
\textbf{Obra} & SOBRENOME, Nome. ``Título do artigo.'' \textit{Nome do
               Periódico}, v. X, n. Y, p. Z--W, Ano. \\
\textbf{Palavras-chave} & Tema A; Abordagem B; Contexto C \\
\midrule
\textbf{Citações diretas} &
  ``Citação direta do artigo (Autor, ano, p. X).'' \\
\midrule
\textbf{Paráfrases} &
  Síntese do argumento central do artigo (Autor, ano, p. X--Y). \\
\midrule
\textbf{Comentário crítico} &
  \textit{Contribuição:} este artigo oferece uma metodologia robusta
  para análise de [fenômeno] em [contexto]. \\
&
  \textit{Limitação:} a amostra restrita a [descrição] limita a
  generalização dos resultados. \\
\midrule
\textbf{Conexões} &
  \begin{itemize}[nosep]
    \item Base metodológica para o Capítulo 3
    \item Resultados contradizem \cite{autor2022}
  \end{itemize} \\
\bottomrule
\end{longtable}

\newpage

\section*{Índice de Fichamentos}

\begin{longtable}{lll}
\toprule
\textbf{\#} & \textbf{Autor (Ano)} & \textbf{Tema Central} \\
\midrule
1 & Sobrenome (2023) & Conceito principal \\
2 & Sobrenome (2021) & Metodologia de análise \\
3 & Sobrenome (2020) & Revisão teórica \\
\bottomrule
\end{longtable}

\section*{Referências}

\bibliographystyle{abntex2-alf}
\bibliography{fichamento_bib}

\end{document}
